% Generated draft: requires mathematical and source review.
\documentclass[11pt]{article}
\usepackage[margin=25mm]{geometry}
\usepackage{fontspec}
\setmainfont{Latin Modern Roman}
\usepackage{amsmath,amssymb,mathtools,unicode-math}
\setmathfont{Latin Modern Math}
\usepackage{enumitem,needspace,xurl,hyperref}
\setlength{\parindent}{0pt}
\setlength{\parskip}{5pt}
\setlength{\emergencystretch}{3em}
\newcommand{\N}{\mathbb{N}}
\newcommand{\Z}{\mathbb{Z}}
\newcommand{\Q}{\mathbb{Q}}
\newcommand{\R}{\mathbb{R}}
\newcommand{\C}{\mathbb{C}}
\newcommand{\F}{\mathbb{F}}
\newcommand{\A}{\mathbb{A}}
\newcommand{\PP}{\mathbb{P}}
\newcommand{\E}{\mathbb{E}}
\newcommand{\eps}{\varepsilon}
\newcommand{\dd}{\,\mathrm{d}}
\begin{document}
\section*{UnsolvedMath ID: 30003655}
\textbf{Source problem code:} catalogue entry
\subsection{Definitions and mathematical statement}
For a nonnegative integer \(m\), \(I\Sigma_m\) denotes the usual fragment of first-order Peano arithmetic whose induction scheme is restricted to \(\Sigma_m\) formulas, with parameters permitted. Here \(\Sigma_0\) consists of bounded formulas, and \(\Sigma_m\) denotes the corresponding level of the arithmetical hierarchy.

For context, if \(X\) is a finite set and \(k\) is a positive integer, write \([X]^k\) for the set of its \(k\)-element subsets. A colouring with \(r\) colours is a map \(c:[X]^k\to\{0,\ldots,r-1\}\), where \(r\geq 1\). A subset \(H\subseteq X\) is homogeneous if \(c\) is constant on \([H]^k\). These definitions describe ordinary finite Ramsey terminology; they do not identify the particular proof-theoretic construction intended by the source.

\textbf{Source question, with unresolved qualifications.} Can the proof-theoretic finite Ramsey theorem construction referred to in the supplied record be formulated using a constant number of colours while retaining the induction fragment \(I\Sigma_{n-2}\), described in that record as optimal?

The record does not specify the construction, the range or role of \(n\), the parameters from which the number of colours must be independent, or the precise proof-theoretic property that must be retained. In the ordinary indexing convention above, \(I\Sigma_{n-2}\) requires an integer \(n\geq 2\); the intended range is not supplied. Nor does the record specify whether optimality concerns provability, unprovability, proof length, or another comparison between arithmetic fragments. Consequently, a fully quantified mathematical assertion cannot be recovered from this record alone. The reference to the Paris--Harrington principle in the supplied bibliography does not specify a replacement statement or its additional hypotheses.
\subsection{Short English statement}
The supplied question asks whether a particular finite Ramsey construction can use a fixed number of colours while preserving its stated optimal induction fragment \(I\Sigma_{n-2}\). The construction, parameter dependencies, and meaning of optimality require clarification before a precise formal statement can be given.
\subsection{Sources}
\begin{itemize}
\item Supplied local UnsolvedMath registry record, permanent ID \(30003655\), title: Constant-Color Proof-Theoretic Finite Ramsey Theorems. Exact source problem code: \texttt{catalogue entry}. Collection: Oberwolfach Reports Open Problems. No individual proposer is identified; the supplied proposed year is 2017.
\item Supplied report citation: \emph{Mathematical Logic: Proof Theory, Constructive Mathematics} (2018), Oberwolfach Reports, pp. 3121--3183, DOI: \url{https://doi.org/10.4171/owr/2017/53}. The literature-source list also labels this reference ``Oberwolfach Report 53, proof theory.''
\item Supplied literature reference: Freund, \emph{Proof lengths for instances of the Paris--Harrington principle}, \url{https://doi.org/10.1016/j.apal.2017.01.004}.
\item Literature assessment dated exactly \texttt{2026-08-21}: the supplied curation labels the problem open and states that proof-length results establish the stated lower-fragment phenomenon, while retaining the optimal induction fragment with a constant number of colours remains open. This is supplied dated curation derived from the OWR record, not an independently verified literature assessment. The record explicitly calls for independent verification of time-sensitive status and cited claims.
\item The supplied statement status is \texttt{corrected\_verified}; its accompanying explanation concerns preservation of mathematical content during notation, typography, and sentence-structure repairs. This metadata is not independent verification of the formulation or its mathematical claims. The cited external sources have not been consulted for this draft.
\end{itemize}
\end{document}
